\subsection{Examples of convolvable measures and functions}

In Props. 3 and 4, \(\mathcal{C}'(\mathrm{G})\) and \(\mathcal{M}(\mathrm{G})\) are equipped with the topology of compact convergence in \(\mathcal{C}(\mathrm{G})\) and \(\mathcal{K}(\mathrm{G})\) , respectively.

PROPOSITION 3. --- Assume \(\chi\) continuous. Let \(\mu \in \mathcal{C}'(\mathbf{G})\) , \(f \in \mathcal{C}(\mathbf{X})\) . Then:

\begin{enumerate}
\def\labelenumi{(\roman{enumi})}
\item
  \(\mu\) and \(f\) are convolvable relative to \(\beta\) .
\item
  Formula (3) of No.~1 defines for every \(x \in \mathbf{X}\) a convolution product \(\mu *^{\beta}f\) that is continuous and is none other than the element \(\pmb{\gamma}_{\chi}(\mu)f\) defined by the continuous representation \(\pmb{\gamma}_{\chi}\) of \(\mathbf{G}\) in \(\mathcal{C}(\mathbf{X})\) ; moreover, the mapping \((\mu, f) \mapsto \mu *^{\beta}f\) is hypocontinuous relative to the equicontinuous subsets of \(\mathcal{C}'(\mathbf{G})\) and the compact subsets of \(\mathcal{C}(\mathbf{X})\) .
\item
  If in addition \(f \in \mathcal{K}(\mathrm{X})\) , then the product \(\mu *^{\beta}f\) of (ii) belongs to \(\mathcal{K}(\mathrm{X})\) and the mapping \((\mu, f) \mapsto \mu *^{\beta}f\) is hypocontinuous relative to the equicontinuous subsets of \(\mathcal{C}'(\mathrm{G})\) and the compact subsets of \(\mathcal{K}(\mathrm{X})\) .
\end{enumerate}

We know that \(\mu\) and \(f\) are convolvable (§3, No.~2, Prop. 8 (i)). On the other hand, with the notations of §2, we have

\[
\boldsymbol {\gamma} _ {\chi} (\mu) f = \int \left(\boldsymbol {\gamma} _ {\chi} (s) f\right) d \mu (s) \in \mathcal {C} (\mathrm{X})
\]

since \(\mathcal{C}(\mathbf{X})\) is complete. In particular, for every \(x\in \mathbf{X}\) ,

\[
\left(\boldsymbol {\gamma} _ {\chi} (\mu) f\right) (x) = \int \left(\boldsymbol {\gamma} _ {\chi} (s) f\right) (x) d \mu (s).
\]

This, combined with Prop. 2 (i), and §2, No.~6, proves (ii). Finally, if \(f \in \mathcal{K}(\mathrm{X})\) then \(\mu * (f \cdot \beta)\) has compact support (§3, No.~2, Prop. 9), therefore \(\mu *^{\beta} f \in \mathcal{K}(\mathrm{X})\) . For, let us consider the continuous representation \(U\) of \(\mathbf{G}\) in the completion \(\mathcal{K}(\mathrm{X})^{\wedge}\) obtained by extending by continuity the continuous operators \(\pmb{\gamma}_{\chi}(s)\) in \(\mathcal{K}(\mathbf{X})\) (§2, No.~1, Remark 3). Let S be the support of \(\mu\) . The functions \(\pmb{\gamma}_{\chi}(s)f\) , for \(s \in \mathbf{S}\) , have their support contained in a fixed compact set K. The set \(\mathcal{K}(\mathbf{X}, \mathbf{K})\) is a complete linear subspace of \(\mathcal{K}(\mathbf{X})\) . Therefore \(U(\mu)f \in \mathcal{K}(\mathbf{X})\) . One sees as before that \(U(\mu)f = \mu^{*\beta}f\) , and (iii) again follows from §2, No.~6.

PROPOSITION 4. --- Assume that G operates properly in X and that \(\chi\) is continuous. Let \(\mu \in \mathcal{M}(\mathrm{G})\) and \(f \in \mathcal{K}(\mathrm{X})\) .

\begin{enumerate}
\def\labelenumi{(\roman{enumi})}
\item
  \(\mu\) and \(f\) are convolvable relative to \(\beta\) .
\item
  Formula (3) of No.~1 defines for every \(x \in \mathbf{X}\) a convolution product \(\mu *^{\beta}f\) that is continuous.
\item
  The mapping \((\mu, f) \mapsto \mu *^{\beta}f\) of \(\mathcal{M}(\mathrm{G}) \times \mathcal{K}(\mathrm{X})\) into \(\mathcal{C}(\mathrm{X})\) is hypocontinuous relative to the bounded subsets of \(\mathcal{M}(\mathrm{G})\) and the compact subsets of \(\mathcal{K}(\mathrm{X})\) that are contained in some subspace \(\mathcal{K}(\mathrm{X}, \mathrm{L})\) (where \(\mathrm{L}\) is a variable compact subset of \(\mathrm{X}\) ).
\end{enumerate}

We know that \(\mu\) and \(f\) are convolvable (§3, No.~2, Prop. 8 (ii)), and it is clear that the integrals occurring in (3) exist for every \(x \in \mathbf{X}\) . Let \(\mathbf{K}\) and \(\mathbf{L}\) be two compact subsets of \(\mathbf{X}\) . There exists a compact subset \(\mathbf{H}\) of \(\mathbf{G}\) such that the relations \(x \in \mathbf{K}\) and \(s^{-1}x \in \mathbf{L}\) imply \(s \in \mathbf{H}\) ; let \(\varphi \in \mathcal{H}_{+}(\mathbf{G})\) with \(\varphi(s) = 1\) for \(s \in \mathbf{H}\) . Then, for \(f \in \mathcal{H}(\mathbf{X}, \mathbf{L})\) and \(x \in \mathbf{K}\) ,

\[
\begin{array}{c} \int f (s ^ {- 1} x) \chi (s ^ {- 1}, x) d \mu (s) = \int f (s ^ {- 1} x) \chi (s ^ {- 1}, x) \varphi (s) d \mu (s) \\ = ((\varphi \cdot \mu) * ^ {\beta} f) (x). \end{array}
\]

Consequently \(\int f(s^{-1}x)\chi(s^{-1},x)d\mu(s)\) is a continuous function of \(x\) and defines a convolution product \(\mu * ^ {\beta} f\in \mathcal{C}(\mathrm{X})\) . Moreover, the mapping \(\mu \mapsto \varphi \cdot \mu\) of \(\mathcal{M}(\mathrm{G})\) into \(\mathcal{C}'(\mathrm{G})\) is continuous for the topologies of compact convergence. Prop. 3 (iii) therefore implies that the mapping \((\mu ,f)\mapsto \mu * ^ {\beta} f\) of \(\mathcal{M}(\mathrm{G})\times \mathcal{K}(\mathrm{X},\mathrm{L})\) into \(\mathcal{C}(\mathrm{X})\) is, for every compact subset L of X, hypocontinuous relative to the compact subsets of \(\mathcal{K}(\mathrm{X},\mathrm{L})\) . In particular, the mapping \((\mu ,f)\mapsto \mu * ^ {\beta} f\) of \(\mathcal{M}(\mathrm{G})\times \mathcal{K}(\mathrm{X})\) into \(\mathcal{C}(\mathrm{X})\) is separately continuous. Since \(\mathcal{K}(\mathrm{X})\) is barreled, this mapping is hypocontinuous relative to the bounded subsets of \(\mathcal{M}(\mathrm{G})\) (TVS, III, §5, No.~3, Prop. 6).

Remark 1. --- Under the hypotheses of Prop. 4, the mapping \(\mu \mapsto \mu *^{\beta}f\) of \(\mathcal{M}_{+}(\mathrm{G})\) into \(\mathcal{C}(\mathrm{X})\) is continuous when \(\mathcal{M}_{+}(\mathrm{G})\) is equipped with the vague topology, for every \(f\in \mathcal{K}(\mathrm{X})\) . For, let K be a compact subset of X, S the (compact) support of \(f\) ; since G operates properly in X, the set of \(s\in \mathrm{G}\) for which there exists an \(x\in \mathrm{K}\) with \(s^{-1}x\in \mathrm{S}\) is a compact subset L of G (GT, III, §4, No.~5, Th. 1). Let \(\varepsilon\) be a number \(>0\) , \(\varphi\) a function in \(\mathcal{K}_{+}(\mathrm{G})\) equal to 1 on the compact set L, \(\mu_0\) an element of \(\mathcal{M}_{+}(\mathrm{G})\) ; the set \(\mathrm{W}_0\) of measures \(\mu \in \mathcal{M}_{+}(\mathrm{G})\) such that

\[
\left| \int \varphi (s) d \mu (s) - \int \varphi (s) d \mu_ {0} (s) \right| \leqslant \varepsilon
\]

is a neighborhood of \(\mu_0\) in \(\mathcal{M}_{+}(\mathrm{G})\) . On the other hand, the function \((s,x)\mapsto f(s^{-1}x)\chi (s^{-1},x)\) is uniformly continuous on \(\mathbf{L}\times \mathbf{K}\) , therefore there exists a finite number of points \(x_{i}\in \mathbf{K}\) ( \(1\leqslant i\leqslant n\) ) such that for every \(x\in \mathbf{K}\) , there is an \(i\) for which

\[
\left| f (s ^ {- 1} x) \chi (s ^ {- 1}, x) - f (s ^ {- 1} x _ {i}) \chi (s ^ {- 1}, x _ {i}) \right| \leqslant \varepsilon
\]

for all \(s \in \mathbf{L}\) . Since \(\mu(\mathbf{L}) \leqslant \int \varphi(s) d\mu_0(s) + \varepsilon\) for all \(\mu \in \mathbf{W}_0\) , also

\[
\begin{array}{r l} \Big | \int f (s ^ {- 1} x) \chi (s ^ {- 1}, x) d \mu (s) - \int f (s ^ {- 1} x _ {i}) \chi (s ^ {- 1}, x _ {i}) d \mu (s) \Big | & \\ & \leqslant \varepsilon \Big (\int \varphi (s) d \mu_ {0} (s) + \varepsilon \Big) \end{array}
\]

for every x satisfying the preceding inequality and every \(\mu \in W_{0}\) . Now let W be the neighborhood of \(\mu_{0}\) in \(\mathcal{M}_{+}(G)\) formed by the measures \(\mu \in W_{0}\) such that

\[
\left| \int f (s ^ {- 1} x _ {i}) \chi (s ^ {- 1}, x _ {i}) d \mu (s) - \int f (s ^ {- 1} x _ {i}) \chi (s ^ {- 1}, x _ {i}) d \mu_ {0} (s) \right| \leqslant \varepsilon
\]

for \(1 \leqslant i \leqslant n\) . It is clear that for every measure \(\mu \in \mathbf{W}\) and every \(x \in \mathbf{K}\) ,

\[
\begin{array}{r l} & {\Big | \int f (s ^ {- 1} x) \chi (s ^ {- 1}, x) d \mu (s) - \int f (s ^ {- 1} x) \chi (s ^ {- 1}, x) d \mu_ {0} (s) \Big |} \\ & {\qquad \leqslant \varepsilon \Big (2 \int \varphi (s) d \mu_ {0} (s) + 2 \varepsilon + 1 \Big),} \end{array}
\]

and since \(\varepsilon\) is arbitrary, this proves our assertion.

PROPOSITION 5. --- Assume \(\chi\) a continuous multiplier and each function \(\chi(s,\cdot)\) bounded.

\begin{enumerate}
\def\labelenumi{(\roman{enumi})}
\item
  The function \(s \mapsto \rho(s) = \sup_{x \in X} \chi(s^{-1}, x)\) on \(G\) is lower semi-continuous \(> 0\) and satisfies \(\rho(st) \leqslant \rho(s)\rho(t)\) for all \(s, t\) in \(G\) .
\item
  Let \(\mu \in \mathcal{M}^{\rho}(\mathrm{G})\) and \(f \in \mathrm{L}^{\infty}(\mathrm{X},\beta)\) . Then \(\mu\) and \(f\) are convolvable and \(\mu *^{\beta} f\) is given locally almost everywhere by the formula (3) of No.~1. One has \(\mu *^{\beta} f \in \mathrm{L}^{\infty}(\mathrm{X},\beta)\) , and \(\| \mu *^{\beta} f \|_{\infty} \leqslant \| \mu \|_{\rho} \| f \|_{\infty}\) .
\item
  If, moreover, \(f \in \mathcal{C}^{\infty}(\mathrm{X})\) (resp. \(\overline{\mathcal{K}(\mathrm{X})}\) ), then formula (3) of No.~1 defines for every \(x\) a convolution product \(\mu *^{\beta}f\) that belongs to \(\mathcal{C}^{\infty}(\mathrm{X})\) (resp. \(\overline{\mathcal{K}(\mathrm{X})}\) ).
\item
  If \(f \in \overline{\mathcal{K}(\mathrm{X})}\) , then the convolution product \(\mu *^{\beta}f\) defined by (3) is none other than the element \(\gamma_{\chi}(\mu)f\) defined by the continuous representation \(\gamma_{\chi}\) of \(G\) in \(\overline{\mathcal{K}(\mathrm{X})}\) .
\end{enumerate}

The identity \(\chi(st,x) = \chi(s,tx)\chi(t,x)\) implies at once that \(\rho(st) \leqslant \rho(s)\rho(t)\) . On the other hand, \(\rho\) is lower semi-continuous, being the upper envelope of continuous functions.

Let \(\mu \in \mathcal{M}^{\rho}(\mathrm{G})\) . By Prop. 1 of No.~1, \(\mu\) and 1 are convolvable; Prop. 2 (i) shows that \((|\mu|^{*\beta}1)(x) \leqslant \int_{\mathrm{G}} \rho(s)d|\mu|(s)\) locally \(\beta\) -almost everywhere. Therefore, if \(f\) is \(\beta\) -measurable and \(|f| \leqslant 1\) , then \(\mu\) and \(f\) are convolvable and \(\mathrm{N}_{\infty}(\mu^{*\beta}f) \leqslant \int \rho(s)d|\mu|(s)\) . Moreover, \(\mu^{*\beta}f\) is given locally almost everywhere by formula (3) of No.~1, because condition (iii) of Prop. 2 of No.~1 is satisfied. This implies (ii).

Suppose \(f\) continuous and bounded by 1 in absolute value. It is clear that the integrals occurring in (3) exist for all \(x \in X\) . Let us show that they depend continuously on \(x\) . We can suppose \(\mu \geqslant 0\) . Let \(x_0 \in X\) and \(\varepsilon > 0\) . Let \(K\) be a compact subset of \(G\) such that \(\int_{G - K} \rho(s) d\mu(s) \leqslant \varepsilon\) . There exists a neighborhood \(V\) of \(x_0\) in \(X\) such that \(x \in V\) implies

\[
\left| f \left(s ^ {- 1} x\right) \chi \left(s ^ {- 1}, x\right) - f \left(s ^ {- 1} x _ {0}\right) \chi \left(s ^ {- 1}, x _ {0}\right) \right| \leqslant \varepsilon / \mu (\mathrm{K})
\]

for \(s \in \mathbf{K}\) . Then, for \(x \in \mathbf{V}\) ,

\[
\begin{array}{r l} & {\Big | \int f (s ^ {- 1} x) \chi (s ^ {- 1}, x) d \mu (s) - \int f (s ^ {- 1} x _ {0}) \chi (s ^ {- 1}, x _ {0}) d \mu (s) \Big |} \\ & {\qquad \leqslant 2 \int_ {\mathrm{G-K}} \rho (s) d \mu (s) + \int_ {\mathrm{K}} \frac {\varepsilon}{\mu (\mathrm{K})} d \mu (s) \leqslant 3 \varepsilon ,} \end{array}
\]

whence our assertion. Suppose that in addition \(f \in \overline{\mathcal{H}(\mathrm{X})}\) . Let \(\mathbf{H}\) be a compact subset of \(\mathbf{X}\) such that \(|f(y)| \leqslant \varepsilon\) for \(y \notin \mathbf{H}\) . Let \(x \notin \mathbf{KH}\) . Then \(s^{-1}x \notin \mathbf{H}\) for \(s \in \mathbf{K}\) , therefore

\[
\begin{array}{r l} \Big | \int_ {\mathrm{G}} f (s ^ {- 1} x) \chi (s ^ {- 1}, x) d \mu (s) \Big | & \leqslant \int_ {\mathrm{G} - \mathrm{K}} \rho (s) d \mu (s) + \int_ {\mathrm{K}} \varepsilon \rho (s) d \mu (s) \\ & \leqslant \varepsilon \Big (1 + \int_ {\mathrm{G}} \rho (s) d \mu (s) \Big), \end{array}
\]

which completes the proof of (iii).

Finally, if \(f \in \overline{\mathcal{K}(\mathbf{X})}\) then, since \(\varepsilon_x \in \mathcal{M}^1(\mathbf{X})\) for all \(x \in \mathbf{X}\) , we have

\[
\left(\boldsymbol {\gamma} _ {\chi} (\mu) f\right) (x) = \int \left(\boldsymbol {\gamma} _ {\chi} (s) f\right) (x) d \mu (s),
\]

thus \(\gamma_{\chi}(\mu)f\) is the convolution product \(\mu *^{\beta}f\) defined by (3).

PROPOSITION 6. --- Assume \(\chi\) a continuous multiplier and each function \(\chi(s, \cdot)\) bounded. Let \(\rho(s) = \sup_{x \in X} \chi(s^{-1}, x)\) . Let \(p\) and \(q\) be two conjugate exponents ( \(1 \leqslant p < +\infty\) ). Let \(\mu \in \mathcal{M}^{\rho^{1/q}}(G)\) and \(f \in L^{p}(X, \beta)\) . Then:

\begin{enumerate}
\def\labelenumi{(\roman{enumi})}
\item
  \(\mu\) and \(f\) are convolvable;
\item
  the convolution product \(\mu * ^ {\beta} f\) is given locally \(\beta\) -almost everywhere by the formula (3), and is equal locally \(\beta\) -almost everywhere to a function \(g\in \mathbf{L}^p (\mathbf{X},\beta)\) such that \(\| g\| _p\leqslant \| \mu \|_{\rho^{1 / q}}\| f\| _p;\) (iii) \(g\) is equal to the element \(\gamma_{\chi}(\mu)f\) defined by the continuous representation \(\gamma_{\chi}\) of \(\mathbf{G}\) in \(\mathbf{L}^p (\mathbf{X},\beta)\) .
\end{enumerate}

We have

\[
\int^ {*} \| \gamma_ {\chi} (s) f \| _ {p} d | \mu | (s) \leqslant \left(\int^ {*} \rho (s) ^ {1 / q} d | \mu | (s)\right) \| f \| _ {p} <   + \infty
\]

by §2, No.~5, formula (5). On the other hand, the mapping \(s \mapsto \pmb{\gamma}_{\chi}(s)f\) of G into \(\mathrm{L}^p (\mathbf{X},\beta)\) is continuous (§2, No.~5, Prop. 9). Therefore this mapping is \(\mu\) -integrable. Let

\[
g = \int_ {\mathrm{G}} \left(\boldsymbol {\gamma} _ {\chi} (s) f\right) d \mu (s) \in \mathrm{L} ^ {p} (\mathrm{X}, \beta).
\]

We have \(\| g\| _p\leqslant \left(\int \rho^{1 / q}(s)d|\mu |(s)\right)\| f\| _p\) . Applying the preceding remarks to \(|f|\) , one sees that the mapping \(s\mapsto \varepsilon_s*\mid f|\) of G into \(\mathbf{L}^p (\mathbf{X},\beta)\) is \(\mu\) -integrable, therefore that, for every \(h\in \mathcal{K}(\mathbf{X})\) , the mapping \(s\mapsto \langle h,\varepsilon_s*(|f|\cdot \beta)\rangle\) is \(\mu\) -integrable. Prop. 7 of §1, No.~5 then proves that \(\mu\) and \(f\cdot \beta\) are convolvable. Moreover,

\[
\begin{array}{r l} \int_ {\mathrm{X}} g (x) h (x) d \beta (x) & = \int_ {\mathrm{G}} d \mu (s) \int_ {\mathrm{X}} (\boldsymbol {\gamma} _ {\chi} (s) f) (x) h (x) d \beta (x) \\ & = \int_ {\mathrm{G}} \langle h, \varepsilon_ {s} * (f \cdot \beta) \rangle d \mu (s), \end{array}
\]

and this last integral is equal to \(\langle h, \mu * (f \cdot \beta) \rangle\) by Prop. 7 of §1, No.~5. One therefore sees that g is a convolution product of \(\mu\) and f.~This convolution product is given locally \(\beta\) -almost everywhere by (3), by Prop. 2 and the Remark that follows it.

By an abuse of notation, it is often one of the functions g of the statement that is denoted \(\mu * ^ {\beta} f\) , which permits writing

\[
\left\| \mu * ^ {\beta} f \right\| _ {p} \leqslant \left\| \mu \right\| _ {\rho^ {1 / q}} \| f \| _ {p}.
\]

If X is countable at infinity, this style of notation is, moreover, entirely justified.

COROLLARY. --- Under the hypotheses of Prop. 6, the mapping \((\mu, f) \mapsto \mu *^{\beta} f\) defines on \(\mathrm{L}^{p}(\mathrm{X}, \beta)\) the structure of a left module over \(\mathcal{M}^{\rho^{1/q}}(\mathrm{G})\) ( \(1 \leqslant p \leqslant +\infty\) ).

This follows from Props. 5 and 6 and the associativity of the convolution product.

Remark 2. --- Let X be a locally compact space on which a locally compact group G operates continuously on the right by \((x,s)\mapsto xs\) . Let \(\beta\) be a positive measure on X. Let \(\chi\) be a function \textgreater0 on \(G\times X\) , measurable for every measure on \(G\times X\) , such that \(\boldsymbol{\delta}(s)\boldsymbol{\beta}=\chi(s,\cdot)\cdot\boldsymbol{\beta}\) for every \(s\in G\) . Let f be a locally \(\beta\) -integrable function on X and let \(\mu\) be a measure on G. If \(f\cdot\beta\) and \(\mu\) are convolvable (for the mapping \((x,s)\mapsto xs\) of \(X\times G\) into X), then \((f\cdot\beta)*\mu\) has base \(\beta\) . We then say that f and \(\mu\) are convolvable relative to \(\beta\) ; every density of \((f\cdot\beta)*\mu\) with respect to \(\beta\) is called a convolution product of f and \(\mu\) relative to \(\beta\) , and is denoted \(f * ^ {\beta} \mu\) or simply \(f*\mu\) .

Let \(G^{0}\) be the group opposite to G. By \((s,x)\mapsto xs\) , \(G^{0}\) operates continuously on the left in X. To say that f and \(\mu\) are convolvable in the foregoing sense is equivalent to saying that \(\mu\) and f are convolvable for \(G^{0}\) operating on the left in X; and the convolution products \(f * ^ {\beta} \mu\) are none other than the convolution products \(\mu * ^ {\beta} f\) for \(G^{0}\) operating on the left in X. On the other hand, for \(s\in G^{0}\) one has \(\pmb{\gamma}(s)\beta = \chi(s^{-1},\cdot)\cdot \beta\) . The results of Nos. 1 and 2 may then be translated immediately into results concerning the products \(f * ^ {\beta} \mu\) . In particular:

\begin{enumerate}
\def\labelenumi{\arabic{enumi})}
\tightlist
\item
  If \(s \in \mathbf{G}\) and \(f\) is locally \(\beta\) -integrable, then \(f\) and \(\varepsilon_s\) are convolvable and
\end{enumerate}

\[
(f * \varepsilon_ {s}) (x) = \chi (s ^ {- 1}, x) f (x s ^ {- 1}).\tag{4}
\]

\begin{enumerate}
\def\labelenumi{\arabic{enumi})}
\setcounter{enumi}{1}
\tightlist
\item
  If \(f\) and \(\mu\) are convolvable and if one of the conditions (i), (ii), (iii) of Prop. 2 is fulfilled, then \(f *^{\beta} \mu\) is given locally \(\beta\) -almost everywhere by
\end{enumerate}

\[
(f * ^ {\beta} \mu) (x) = \int_ {G} f (x s ^ {- 1}) \chi (s ^ {- 1}, x) d \mu (s).\tag{5}
\]

We leave to the reader the task of translating the other statements. Note that if \(\chi\) is continuous and \(\beta\) has support X, then

\[
\chi (t s, x) = \chi (s, x t) \chi (t, x) \quad (x \in \mathrm{X}; s, t \text {in} \mathrm{G}).\tag{6}
\]

